\documentclass[../main.tex]{subfiles}

\begin{document}

\setcounter{chapter}{-1}
\chapter{集合、映射与关系的预备知识}\label{chapter:ch00-预备知识}

在预备知识中, 我们描述一些本书始终会用到的集合、映射与关系的概念, 给出的论据倾向于合理与易理解, 而非基于集合论公理的严格证明. 在称为关于集合的 Zermelo-Frankel 公理系统之上, 可以正式地建立集合、关系以及函数的性质. 有兴趣的读者可以查阅 John Kelley 的书《General Topology》\cite{Kel75}、Paul Halmos 的书《Naive Set Theory》\cite{Hal98} 以及 Thomas Jech 的书《Set Theory》\cite{Jec06} 的引言与附录.

\subfile{Chapter-00/section-01.tex}

\subfile{Chapter-00/section-02.tex}

\subfile{Chapter-00/section-03.tex}

\end{document}